\chapter{M\'odulos de base}
\label{cap:base}

Este cap\'itulo documenta os m\'odulos que n\~ao pertencem a um driver
espec\'ifico: \code{main.rs}, \code{arch.rs} e \code{sync.rs}. Eles fornecem o
ponto de entrada, as macros de impress\~ao, o acesso a portas de I/O, o
desligamento/reinicializa\c{c}\~ao e o spinlock do console.

\section{\code{main.rs}: ponto de entrada Rust}

\subsection{Atributos globais}

\begin{lstlisting}[caption={Atributos e m\'odulos declarados em \code{main.rs}.},label={lst:mainattrs}]
#![no_std]
#![no_main]

mod arch;
mod boot;
mod gdt;
mod idt;
mod keyboard;
mod paging;
mod serial;
mod shell;
mod sync;
mod user;
mod vga;
\end{lstlisting}

\begin{itemize}
  \item \code{no\_std}: o crate n\~ao usa a biblioteca padr\~ao; n\~ao h\'a
        \code{std::io}, \code{String}, \code{Vec}, sistema de arquivos nem
        \emph{threads}.
  \item \code{no\_main}: n\~ao existe a fun\c{c}\~ao \code{main} do runtime
        Rust. O ponto de entrada real \'e \code{\_start}, definido no
        assembly de \code{boot.rs}.
  \item \code{mod boot}: necess\'ario para que o linker inclua o c\'odigo de
        montagem e o cabe\c{c}alho Multiboot2.
\end{itemize}

\subsection{\code{kernel\_main}}

\begin{lstlisting}[caption={Assinatura e corpo principal de \code{kernel\_main}.},label={lst:kernelmain}]
#[unsafe(no_mangle)]
pub extern "C" fn kernel_main(multiboot_magic: u32, multiboot_info: u32) -> ! {
    vga::init();
    serial::init();
    paging::init(multiboot_magic, multiboot_info);
    gdt::init();
    keyboard::init();
    idt::init();
    arch::enable_interrupts();

    println!();
    println!("==============================================================");
    println!("  Rust VGA Kernel - x86_64 / modo longo");
    /* ... banner ... */

    if multiboot_magic != 0x36d7_6289 {
        println!("AVISO: assinatura multiboot2 invalida: {:#010x}",
                 multiboot_magic);
    }
    println!("Multiboot2: magic={:#010x} info={:#010x}",
             multiboot_magic, multiboot_info);
    println!();
    println!("Digite 'help' para ver os comandos.");
    println!();

    shell::run()
}
\end{lstlisting}

\begin{table}[H]
\centering
\caption{Contrato de \code{kernel\_main}.}
\label{tab:kernelmain}
\begin{tabularx}{\linewidth}{@{}lX@{}}
\toprule
\textbf{Aspecto} & \textbf{Descri\c{c}\~ao} \\
\midrule
Visibilidade & \code{pub extern "C"}, exportada sem \emph{mangling}. \\
Argumentos & \code{multiboot\_magic} (deve ser \code{0x36D76289}) e
\code{multiboot\_info} (endere\c{c}o f\'isico da estrutura Multiboot2). \\
Retorno & \code{!}: nunca retorna; termina em \code{shell::run()}. \\
Pr\'e-condi\c{c}\~oes & CPU em modo longo, pagina\c{c}\~ao ativa, pilha alta
v\'alida, \code{.bss} zerada pelo carregador. \\
Efeitos & Inicializa VGA, serial, alocador de frames, GDT/TSS, teclado,
IDT/PIC/PIT e liga interrup\c{c}\~oes; imprime banner; entra no shell. \\
Chamado por & Stub de assembly em \code{boot.rs}. \\
\bottomrule
\end{tabularx}
\end{table}

O valor \code{multiboot\_info} \'e usado por \code{paging::init} para
percorrer as tags Multiboot2 e liberar os frames de RAM dispon\'ivel. A
estrutura tamb\'em \'e reservada para n\~ao ser entregue pelo alocador.

\subsection{Macros \code{print!} e \code{println!}}

As macros s\~ao definidas com \code{macro\_rules!} e \code{\$crate}, para
funcionar em qualquer m\'odulo do crate sem importa\c{c}\~ao:

\begin{lstlisting}[caption={Macros de impress\~ao sem aloca\c{c}\~ao.},label={lst:macros}]
#[macro_export]
macro_rules! print {
    ($($arg:tt)*) => {
        $crate::tty::_print(format_args!($($arg)*))
    };
}

#[macro_export]
macro_rules! println {
    () => { $crate::print!("\n") };
    ($($arg:tt)*) => {
        $crate::print!("{}\n", format_args!($($arg)*))
    };
}
\end{lstlisting}

\code{format\_args!} cria uma descri\c{c}\~ao de formata\c{c}\~ao sem alocar;
a implementa\c{c}\~ao de \code{fmt::Write} em \code{tty.rs} consome os
fragmentos, anexa ao buffer do console ativo e espelha na serial.
N\~ao h\'a \code{format!} (que exigiria \code{alloc}).

\subsection{Panic handler}

\begin{lstlisting}[caption={\emph{Panic handler} do kernel.},label={lst:panic}]
#[panic_handler]
fn panic(info: &PanicInfo) -> ! {
    crate::println!();
    crate::println!("[PANIC] {}", info);
    crate::arch::halt_loop()
}
\end{lstlisting}

Como o perfil usa \code{panic = "abort"}, o handler n\~ao precisa de
\emph{unwinding}. Ele imprime a mensagem no console e para a CPU. Se o
p\^anico ocorrer antes de \code{vga::init()}, a sa\'ida pode n\~ao ser
vis\'ivel; se ocorrer dentro do console com o spinlock adquirido, pode
ocorrer \emph{deadlock} (o handler tentaria imprimir e esperaria o lock).

\section{\code{arch.rs}: arquitetura e portas de I/O}

\subsection{Portas de I/O}

O m\'odulo \code{arch::port} encapsula as instru\c{c}\~oes \code{in} e
\code{out} do x86:

\begin{lstlisting}[caption={Leitura e escrita de portas de 8 bits.},label={lst:portio}]
pub unsafe fn inb(port: u16) -> u8 {
    let value: u8;
    asm!("in al, dx",
         in("dx") port,
         out("al") value,
         options(nomem, nostack, preserves_flags));
    value
}

pub unsafe fn outb(port: u16, value: u8) {
    asm!("out dx, al",
         in("dx") port,
         in("al") value,
         options(nomem, nostack, preserves_flags));
}
\end{lstlisting}

As fun\c{c}\~oes s\~ao \code{unsafe} porque qualquer porta pode n\~ao existir
ou pode ter efeitos colaterais no hardware. As op\c{c}\~oes \code{nomem},
\code{nostack} e \code{preserves\_flags} informam ao compilador que a
instru\c{c}\~ao n\~ao altera mem\'oria, pilha ou flags, permitindo otimiza\c{c}\~oes
e evitando que o bloco seja removido.

\begin{table}[H]
\centering
\caption{API de \code{arch}.}
\label{tab:archapi}
\begin{tabularx}{\linewidth}{@{}lX@{}}
\toprule
\textbf{Fun\c{c}\~ao} & \textbf{Descri\c{c}\~ao} \\
\midrule
\code{inb(port)} & L\^e um byte de uma porta de I/O. \code{unsafe}. \\
\code{outb(port, value)} & Escreve um byte em uma porta. \code{unsafe}. \\
\code{halt\_loop()} & Executa \code{cli; hlt} em la\c{c}o infinito; n\~ao
retorna. Usado no panic e no comando \code{halt}. \\
\code{reboot()} & Espera o buffer de entrada do 8042 esvaziar e envia o
comando \code{0xFE} (reset). Se n\~ao funcionar, cai em \code{halt\_loop}. \\
\bottomrule
\end{tabularx}
\end{table}

\subsection{Reinicializa\c{c}\~ao}

O reset usa o controlador 8042, que historicamente controla a linha de reset
da CPU. O c\'odigo espera at\'e 100.000 itera\c{c}\~oes pelo bit
\code{INPUT\_FULL} (\code{0x02}) do status limpar e escreve \code{0xFE} na
porta de comando (\code{0x64}). Em m\'aquinas virtualizadas ou sem 8042, o
reset pode n\~ao ocorrer; nesse caso a CPU \'e parada.

\section{\code{sync.rs}: spinlock}

\subsection{Implementa\c{c}\~ao}

\begin{lstlisting}[caption={Spinlock com \code{AtomicBool}.},label={lst:spinlock}]
pub struct SpinLock<T> {
    locked: AtomicBool,
    data: UnsafeCell<T>,
}

unsafe impl<T: Send> Sync for SpinLock<T> {}
unsafe impl<T: Send> Send for SpinLock<T> {}

impl<T> SpinLock<T> {
    pub const fn new(data: T) -> Self { /* ... */ }

    pub fn lock(&self) -> SpinLockGuard<'_, T> {
        while self.locked.swap(true, Ordering::Acquire) {
            core::hint::spin_loop();
        }
        SpinLockGuard { lock: self }
    }
}
\end{lstlisting}

O \emph{guard} implementa \code{Deref}, \code{DerefMut} e \code{Drop}. Ao ser
destru\'ido, armazena \code{false} com ordena\c{c}\~ao \code{Release}. O
spinlock \'e usado apenas para o console:

\begin{lstlisting}[caption={Console global protegido por spinlock.},label={lst:consoleglobal}]
static CONSOLE: SpinLock<Console> = SpinLock::new(Console::new());
\end{lstlisting}

Como \code{Console::new()} \'e \code{const}, o lock pode ser inicializado
estaticamente sem c\'odigo de inicializa\c{c}\~ao em tempo de execu\c{c}\~ao.

\subsection{Propriedades e limita\c{c}\~oes}

\begin{itemize}
  \item \textbf{Bloqueio girat\'orio}: se outra CPU ou interrup\c{c}\~ao
        mantiver o lock, a CPU atual gira at\'e liberar. N\~ao h\'a
        \emph{sleep} nem timeout.
  \item \textbf{Sem IRQ save/restore}: desabilitar interrup\c{c}\~oes ao
        adquirir o lock n\~ao \'e feito. Com IRQs desligadas no boot, isso
        n\~ao \'e problema hoje; ao habilitar IRQs, um handler que use o
        console pode causar \emph{deadlock}.
  \item \textbf{Sem detec\c{c}\~ao de recurs\~ao}: adquirir o mesmo lock duas
        vezes na mesma CPU trava.
  \item \textbf{Ordem de mem\'oria}: \code{Acquire}/\code{Release} garante
        que escritas ao console n\~ao sejam reordenadas para fora da
        regi\~ao cr\'itica em arquiteturas com mem\'oria fraca; no x86,
        o modelo \'e mais forte, mas a anota\c{c}\~ao mant\'em o c\'odigo
        correto.
\end{itemize}

\begin{nota}[Escopo atual]
Com uma CPU, interrup\c{c}\~oes desligadas e apenas o shell chamando
\code{print!}, o spinlock nunca fica contended. Ele existe como prepara\c{c}\~ao
para extens\~oes e para documentar a inten\c{c}\~ao de exclus\~ao m\'utua.
\end{nota}
